% Options for packages loaded elsewhere
% Options for packages loaded elsewhere
\PassOptionsToPackage{unicode}{hyperref}
\PassOptionsToPackage{hyphens}{url}
\PassOptionsToPackage{dvipsnames,svgnames,x11names}{xcolor}
%
\documentclass[
  11pt,
]{article}
\usepackage{xcolor}
\usepackage[margin=1in]{geometry}
\usepackage{amsmath,amssymb}
\setcounter{secnumdepth}{-\maxdimen} % remove section numbering
\usepackage{iftex}
\ifPDFTeX
  \usepackage[T1]{fontenc}
  \usepackage[utf8]{inputenc}
  \usepackage{textcomp} % provide euro and other symbols
\else % if luatex or xetex
  \usepackage{unicode-math} % this also loads fontspec
  \defaultfontfeatures{Scale=MatchLowercase}
  \defaultfontfeatures[\rmfamily]{Ligatures=TeX,Scale=1}
\fi
\usepackage{lmodern}
\ifPDFTeX\else
  % xetex/luatex font selection
\fi
% Use upquote if available, for straight quotes in verbatim environments
\IfFileExists{upquote.sty}{\usepackage{upquote}}{}
\IfFileExists{microtype.sty}{% use microtype if available
  \usepackage[]{microtype}
  \UseMicrotypeSet[protrusion]{basicmath} % disable protrusion for tt fonts
}{}
\makeatletter
\@ifundefined{KOMAClassName}{% if non-KOMA class
  \IfFileExists{parskip.sty}{%
    \usepackage{parskip}
  }{% else
    \setlength{\parindent}{0pt}
    \setlength{\parskip}{6pt plus 2pt minus 1pt}}
}{% if KOMA class
  \KOMAoptions{parskip=half}}
\makeatother
% Make \paragraph and \subparagraph free-standing
\makeatletter
\ifx\paragraph\undefined\else
  \let\oldparagraph\paragraph
  \renewcommand{\paragraph}{
    \@ifstar
      \xxxParagraphStar
      \xxxParagraphNoStar
  }
  \newcommand{\xxxParagraphStar}[1]{\oldparagraph*{#1}\mbox{}}
  \newcommand{\xxxParagraphNoStar}[1]{\oldparagraph{#1}\mbox{}}
\fi
\ifx\subparagraph\undefined\else
  \let\oldsubparagraph\subparagraph
  \renewcommand{\subparagraph}{
    \@ifstar
      \xxxSubParagraphStar
      \xxxSubParagraphNoStar
  }
  \newcommand{\xxxSubParagraphStar}[1]{\oldsubparagraph*{#1}\mbox{}}
  \newcommand{\xxxSubParagraphNoStar}[1]{\oldsubparagraph{#1}\mbox{}}
\fi
\makeatother


\usepackage{longtable,booktabs,array}
\newcounter{none} % for unnumbered tables
\usepackage{calc} % for calculating minipage widths
% Correct order of tables after \paragraph or \subparagraph
\usepackage{etoolbox}
\makeatletter
\patchcmd\longtable{\par}{\if@noskipsec\mbox{}\fi\par}{}{}
\makeatother
% Allow footnotes in longtable head/foot
\IfFileExists{footnotehyper.sty}{\usepackage{footnotehyper}}{\usepackage{footnote}}
\makesavenoteenv{longtable}
\usepackage{graphicx}
\makeatletter
\newsavebox\pandoc@box
\newcommand*\pandocbounded[1]{% scales image to fit in text height/width
  \sbox\pandoc@box{#1}%
  \Gscale@div\@tempa{\textheight}{\dimexpr\ht\pandoc@box+\dp\pandoc@box\relax}%
  \Gscale@div\@tempb{\linewidth}{\wd\pandoc@box}%
  \ifdim\@tempb\p@<\@tempa\p@\let\@tempa\@tempb\fi% select the smaller of both
  \ifdim\@tempa\p@<\p@\scalebox{\@tempa}{\usebox\pandoc@box}%
  \else\usebox{\pandoc@box}%
  \fi%
}
% Set default figure placement to htbp
\def\fps@figure{htbp}
\makeatother


% definitions for citeproc citations
\NewDocumentCommand\citeproctext{}{}
\NewDocumentCommand\citeproc{mm}{%
  \begingroup\def\citeproctext{#2}\cite{#1}\endgroup}
\makeatletter
 % allow citations to break across lines
 \let\@cite@ofmt\@firstofone
 % avoid brackets around text for \cite:
 \def\@biblabel#1{}
 \def\@cite#1#2{{#1\if@tempswa , #2\fi}}
\makeatother
\newlength{\cslhangindent}
\setlength{\cslhangindent}{1.5em}
\newlength{\csllabelwidth}
\setlength{\csllabelwidth}{3em}
\newenvironment{CSLReferences}[2] % #1 hanging-indent, #2 entry-spacing
 {\begin{list}{}{%
  \setlength{\itemindent}{0pt}
  \setlength{\leftmargin}{0pt}
  \setlength{\parsep}{0pt}
  % turn on hanging indent if param 1 is 1
  \ifodd #1
   \setlength{\leftmargin}{\cslhangindent}
   \setlength{\itemindent}{-1\cslhangindent}
  \fi
  % set entry spacing
  \setlength{\itemsep}{#2\baselineskip}}}
 {\end{list}}
\usepackage{calc}
\newcommand{\CSLBlock}[1]{\hfill\break\parbox[t]{\linewidth}{\strut\ignorespaces#1\strut}}
\newcommand{\CSLLeftMargin}[1]{\parbox[t]{\csllabelwidth}{\strut#1\strut}}
\newcommand{\CSLRightInline}[1]{\parbox[t]{\linewidth - \csllabelwidth}{\strut#1\strut}}
\newcommand{\CSLIndent}[1]{\hspace{\cslhangindent}#1}



\setlength{\emergencystretch}{3em} % prevent overfull lines

\providecommand{\tightlist}{%
  \setlength{\itemsep}{0pt}\setlength{\parskip}{0pt}}



 


% tinytable LaTeX dependencies. See supplement.qmd for rationale.
\usepackage{tabularray}
\usepackage{float}
\usepackage{graphicx}
\usepackage[normalem]{ulem}
\providecommand{\tinytableTabularrayUnderline}[1]{\underline{#1}}
\providecommand{\tinytableTabularrayStrikeout}[1]{\sout{#1}}
% Author block: Quarto's default \author{names \and …} gives
% names only. We override with authblk to pin affiliations from
% YAML to the title page.
\usepackage[affil-it]{authblk}
\renewcommand\Authsep{, }
\renewcommand\Authand{, }
\renewcommand\Authands{, }
\renewcommand\Affilfont{\itshape\small}
\let\OldAuthor\author
\renewcommand{\author}[1]{}
\AtBeginDocument{%
  \OldAuthor[1]{Daniel W.A. Noble\thanks{Corresponding author: \texttt{daniel.noble@anu.edu.au}}}%
  \OldAuthor[1]{Pieter A. Arnold}%
  \OldAuthor[2,3]{Patrice Pottier}%
  \affil[1]{Division of Ecology and Evolution, Research School of Biology, Australian National University, Canberra, Australia}%
  \affil[2]{Department of Biological and Environmental Sciences, University of Gothenburg, Gothenburg, Sweden}%
  \affil[3]{Evolution \& Ecology Research Centre, School of Biological, Earth and Environmental Sciences, University of New South Wales, Sydney, New South Wales, Australia}%
}
\makeatletter
\@ifpackageloaded{caption}{}{\usepackage{caption}}
\AtBeginDocument{%
\ifdefined\contentsname
  \renewcommand*\contentsname{Table of contents}
\else
  \newcommand\contentsname{Table of contents}
\fi
\ifdefined\listfigurename
  \renewcommand*\listfigurename{List of Figures}
\else
  \newcommand\listfigurename{List of Figures}
\fi
\ifdefined\listtablename
  \renewcommand*\listtablename{List of Tables}
\else
  \newcommand\listtablename{List of Tables}
\fi
\ifdefined\figurename
  \renewcommand*\figurename{Figure}
\else
  \newcommand\figurename{Figure}
\fi
\ifdefined\tablename
  \renewcommand*\tablename{Table}
\else
  \newcommand\tablename{Table}
\fi
}
\@ifpackageloaded{float}{}{\usepackage{float}}
\floatstyle{ruled}
\@ifundefined{c@chapter}{\newfloat{codelisting}{h}{lop}}{\newfloat{codelisting}{h}{lop}[chapter]}
\floatname{codelisting}{Listing}
\newcommand*\listoflistings{\listof{codelisting}{List of Listings}}
\makeatother
\makeatletter
\makeatother
\makeatletter
\@ifpackageloaded{caption}{}{\usepackage{caption}}
\@ifpackageloaded{subcaption}{}{\usepackage{subcaption}}
\makeatother
\usepackage{bookmark}
\IfFileExists{xurl.sty}{\usepackage{xurl}}{} % add URL line breaks if available
\urlstyle{same}
\hypersetup{
  pdftitle={A flexible modelling framework for estimating thermal tolerance and sensitivity},
  pdfkeywords={thermal death time, tolerance landscape, hierarchical
Bayesian model, 4-parameter logistic, heat injury, CT\_max, ectotherm},
  colorlinks=true,
  linkcolor={blue},
  filecolor={Maroon},
  citecolor={Blue},
  urlcolor={Blue},
  pdfcreator={LaTeX via pandoc}}


\title{A flexible modelling framework for estimating thermal tolerance
and sensitivity}
\author{Daniel W.A. Noble \and Pieter A. Arnold \and Patrice Pottier}
\date{}
\begin{document}
\maketitle
\begin{abstract}
The thermal load sensitivity (TLS) and thermal death time (TDT)
frameworks are gold standard approaches for estimating thermal tolerance
and predict organismal vulnerability under natural thermal regimes.
However, these frameworks are usually implemented through a two-step
pipeline. Median times to lethal or sublethal endpoints (\(t_{50}\)) are
first estimated independently across a range of assay temperatures, and
these \(\log_{10}(t_{50})\) estimates are then regressed against
temperature. This two-step approach violates key statistical assumptions
that are likely to bias the mean and variance in derived quantities such
as thermal sensitivity (\(z\)), critical thermal limits (\(CT_{max}\)),
heat-injury predictions, and survival probabilities.

We present a flexible Bayesian four-parameter logistic (4PL) model that
estimates survival or sublethal responses as a joint function of
temperature and exposure duration. We demonstrate how such an approach
recovers all the classical outputs of the TLS and TDT two-stage
modelling frameworks (i.e., \(z\), \(CT_{max}\), heat-injury integral,
survival probabilities), while propagating uncertainty through all
derived quantities. Using simulated data and empirical case studies in
brown shrimp (\emph{Crangon crangon}), zebrafish (\emph{Danio rerio}),
and XXX (* XXX*), we demonstrate the benefits of this new modelling
framework.

We show that the classical and Bayesian frameworks are equivalent, but
that the Bayesian approach delivers properly propagated uncertainty, can
account for complex hierarchical data (i.e., random effects), can
accommodate censoring and overdispersion, and unlocks outputs the
two-step pipeline cannot provide, including full survival surfaces,
mortality percentiles, decomposable between-group contrasts, and
variation in the shape of thermal dose-response curves. We provide a new
R package (\(bayesTLS\)) and recommend the bayesian four-parameter
loigistic model as the default approach for analysing TDT and TLS data,
particularly when the goal is to make accurate predictions of survival
under natural thermal regimes.
\end{abstract}


\section{Introduction}\label{introduction}

Climate change is driving more frequent and extreme heat events
globally, compromising organism fitness
(\citeproc{ref-arnold_thermal_2025}{Arnold \emph{et al.} 2025};
\citeproc{ref-rezende_drosophila_mortality_2020}{Rezende \emph{et al.}
2020}; \citeproc{ref-snook_fertility_2026}{Snook \emph{et al.} 2026}),
with cascading consequences for population persistence and biodiversity
(\citeproc{ref-smale_marine_2019}{Smale \emph{et al.} 2019};
\citeproc{ref-wiens_localextinctions_2016}{Wiens 2016}). Quantifying
differences in thermal tolerance across species is now central to
forecasting climate change impacts on organisms in nature
(\citeproc{ref-deutsch_latitude_2008}{Deutsch \emph{et al.} 2008};
\citeproc{ref-jorgensen_nature_2022}{Jørgensen \emph{et al.} 2022};
\citeproc{ref-pinsky_marine_2019}{Pinsky \emph{et al.} 2019};
\citeproc{ref-sunday_global_2011}{Sunday \emph{et al.} 2011},
\citeproc{ref-sunday_redistribution_2012}{2012}). Threshold-based
measures of heat tolerance, such as \(CT_{max}\), have long been used to
establish differences in climate vulnerability among taxa and predicting
climate change risk (\citeproc{ref-deutsch_latitude_2008}{Deutsch
\emph{et al.} 2008}; \citeproc{ref-pinsky_marine_2019}{Pinsky \emph{et
al.} 2019}; \citeproc{ref-pottier_amphibians_2025}{Pottier \emph{et al.}
2025}; \citeproc{ref-sunday_global_2011}{Sunday \emph{et al.} 2011},
\citeproc{ref-sunday_redistribution_2012}{2012}). While useful,
threshold-based measures do not capture how the negative effects of
extreme heat on fitness depend on both the duration and intensity of
heat stress, limiting our ability to make realistic predictions under
fluctuating temperature regimes that occur in nature
(\citeproc{ref-jorgensen_unifying_2021}{Jørgensen \emph{et al.} 2021};
\citeproc{ref-orsted_finding_2022}{Ørsted \emph{et al.} 2022};
\citeproc{ref-rezende_landscapes_2014}{Rezende \emph{et al.} 2014}).

The Thermal Load Sensitivity (TLS) framework (or equivalently Thermal
Death Time framework) (\citeproc{ref-arnold_thermal_2025}{Arnold
\emph{et al.} 2025}; \citeproc{ref-jorgensen_drosophila_2019}{Jørgensen
\emph{et al.} 2019}, \citeproc{ref-jorgensen_nature_2022}{2022};
\citeproc{ref-noble2026}{Noble \emph{et al.} 2026};
\citeproc{ref-orsted_finding_2022}{Ørsted \emph{et al.} 2022};
\citeproc{ref-rezende_landscapes_2014}{Rezende \emph{et al.} 2014})
resolves these limitations by jointly modelling the effects of heat
intensity and duration on fitness (e.g., survival, fertility,
reproduction). This matters because physiological damage from heat
stress accumulates in a dose-dependent manner --- physiological damage
accumulates faster at higher temperatures and over longer exposures
(\citeproc{ref-arnold_thermal_2025}{Arnold \emph{et al.} 2025};
\citeproc{ref-jorgensen_drosophila_2019}{Jørgensen \emph{et al.} 2019};
\citeproc{ref-orsted_finding_2022}{Ørsted \emph{et al.} 2022};
\citeproc{ref-rezende_landscapes_2014}{Rezende \emph{et al.} 2014}).
Organisms experiencing fluctuating thermal environments in nature can
therefore accumulate damage to physiological systems that results in
sublethal injury that compounds across successive heat exposures if not
fully repaired (\citeproc{ref-arnold_thermal_2025}{Arnold \emph{et al.}
2025}; \citeproc{ref-buckley_damage_2025}{Buckley \emph{et al.} 2025};
\citeproc{ref-jorgensen_unifying_2021}{Jørgensen \emph{et al.} 2021}).
Such processes can lead to delayed responses to extreme heat events.

Thermal load sensitivity modelling, while powerful, traditionally
involves a two-step modelling approach
(\citeproc{ref-jorgensen_drosophila_2019}{Jørgensen \emph{et al.} 2019};
\citeproc{ref-li2023}{Li \emph{et al.} 2023};
\citeproc{ref-orsted_finding_2022}{Ørsted \emph{et al.} 2022};
\citeproc{ref-rezende_drosophila_mortality_2020}{Rezende \emph{et al.}
2020}; \citeproc{ref-truebano2018}{Truebano \emph{et al.} 2018}). First,
experiments will typically be expose organisms to a range of
temperatures for various durations of exposure. The proportion of
live/dead, fertile/infertile, or live/knockdown, for instance, is then
modelled as a function of log transformed exposure time (in minutes or
hours) using dose-response curves fit across all temperatures (usually
at each temperature treatment independently). The time to death,
fertility, or heat knockdown is then derived from these curves, based on
the temperature predicted to reduce these responses by a given threshold
(e.g., \(t_{50}\), the time to reach 50\% mortality). Second, the log10
transformed \(t_{50}\) and temperature are regressed against each other
using a linear model. Thermal sensitivity (i.e., \emph{z}, the inverse
of the slope of this regression) and static critical thermal maxima
(i.e., the intercept of the regression, \(sCT_{max_{time}}\)) are then
derived from this linear model
(\citeproc{ref-orsted_finding_2022}{Ørsted \emph{et al.} 2022};
\citeproc{ref-rezende_landscapes_2014}{Rezende \emph{et al.} 2014}).
\(sCT_{max_{time}}\) is usually computed for a 1 min
(\citeproc{ref-rezende_landscapes_2014}{Rezende \emph{et al.} 2014}), 1
hour (\citeproc{ref-jorgensen_drosophila_2019}{Jørgensen \emph{et al.}
2019}), or 4 hour (\citeproc{ref-orsted_suzukii_2024}{Ørsted \emph{et
al.} 2024}) duration of exposure to heat stress. The 1 hour reference
falls within typical assay durations avoiding substantial extrapolation
which is why it is often recommended
(\citeproc{ref-jorgensen_drosophila_2019}{Jørgensen \emph{et al.}
2019}). Thermal sensitivity and \(CT_{max_{1hour}}\) are than used to
quantify Heat Injury (HI) (e.g., Jørgensen \emph{et al.}
(\citeproc{ref-jorgensen_drosophila_2019}{2019}); Villeneuve \& White
(\citeproc{ref-villeneuve_dynamic_2024}{2024})) in response to realistic
temperatures in nature. Heat injury measures the cumulative percentage
or probability until reaching a given threshold (e.g., \(t_{50}\)), at
which point we would predict 50\% of the population to have died. Recent
work has also emphasized the importance of repair in heat injury
dynamics (\citeproc{ref-arnold_thermal_2025}{Arnold \emph{et al.} 2025};
\citeproc{ref-buckley_damage_2025}{Buckley \emph{et al.} 2025}), and new
modelling approaches have recently been developed to capture these
damage-repair dynamics (\citeproc{ref-arnold_thermal_2025}{Arnold
\emph{et al.} 2025}).

While the two-step approach to modelling thermal sensitivity is
intuitive and simple in practice, it has numerous practical and
statistical limitations. First, the two-step pipeline does not typically
allow correct uncertainty propagation. The \(t_{50}\) point estimates
have non-negligible sampling variance. Treating them as fixed inputs to
the regression discards that sampling variance, which can lead to
miscalibrated standard errors on the slope (\(z\)) and intercept
(\(CT_{max}\)), and inflated \(R^2\) -- a known issue in the statistical
literature (\citeproc{ref-burke_ipd_2017}{Burke \emph{et al.} 2017};
\citeproc{ref-carroll_measurement_2006}{Carroll \emph{et al.} 2006}).
Second, group comparisons (between species, sexes, or life stages)
inherit this miscalibrated uncertainty. As such, differences in \(z\)
between groups cannot be attributed to differences in the upper
asymptote, lower asymptote, or curve steepness of the dose-response
curve, because all three are collapsed into a single \(t_{50}\) without
sampling variance before the comparison is made. Third, the
\(t_{50}\)-extraction step collapses the within-replicate count, and
prevents the hierarchical decomposition of variance across replicates,
plates, populations, sexes, and life stages that could affect parameters
of dose-response curves (\citeproc{ref-harrison_olre_2014}{Harrison
2014}; \citeproc{ref-harrison_mixed_2018}{Harrison \emph{et al.} 2018};
\citeproc{ref-nakagawa_santos_2012}{Nakagawa \& Santos 2012}). Fourth,
the intrinsically hierarchical structure of the experimental data(e.g.,
animals within batches and days of experiments) is largely ignored,
which can lead to pseudoreplicaton, resulting in anticonservative
standard errors that affect biological inference. Finally, propagating
uncertainty from a two-step fit into downstream predictions --- such as
cumulative heat injury --- is rarely attempted in practice understating
the credible range of extreme heat impacts in nature.

Here, we present an alternative to the two-step pipeline: a single
hierarchical four-parameter logistic (4PL) model fit jointly across all
temperature by exposure duration data. The 4PL is a flexible sigmoidal
form that pairs with different sampling distributions --- binomial for
counts of dead or infertile individuals, time-to-event for latencies
(e.g., time to knockdown), Beta for continuous proportions, or Gaussian
for continuous performance measures
(\citeproc{ref-nielsen_nlmm_2004}{Nielsen \emph{et al.} 2004};
\citeproc{ref-ritz_ecotox_2010}{Ritz 2010};
\citeproc{ref-ritz_drc_2015}{Ritz \emph{et al.} 2015};
\citeproc{ref-ritz_streibig_2008}{Ritz \& Streibig 2008};
\citeproc{ref-rudemo_random_1989}{Rudemo \emph{et al.} 1989}). We focus
on binomial outcomes for clarity, since these dominate the TLS
literature (\citeproc{ref-orsted_finding_2022}{Ørsted \emph{et al.}
2022}, \citeproc{ref-orsted_suzukii_2024}{2024};
\citeproc{ref-rezende_drosophila_mortality_2020}{Rezende \emph{et al.}
2020}), but the derivations below apply unchanged when the likelihood is
swapped. Additionally, while we use the four-parameter log-logistic form
throughout for its flexibility other monotonic sigmoidal families (e.g.,
Weibull, log-normal) and simpler nested versions of the log-logistic
(3PL, with one asymptote fixed; 2PL, with both fixed at 0 and 1) can be
used within this framework (\citeproc{ref-ritz_drc_2015}{Ritz \emph{et
al.} 2015}; \citeproc{ref-ritz_streibig_2008}{Ritz \& Streibig 2008}).
We show how all the classical thermal load sensitivity parameters
(\emph{z}, \(CT_{max}\)) can be derived from this single model, while
also easily allowing for predictions of heat injury accumulation,
survival, and their uncertainty in fluctuating thermal regimes. We
develop these models within a Bayesian fr amework which allows for
better uncertainty calibration and error propagation with complex
hierarchical data structures.

\section{A non-linear Bayesian hierarchical model for characterizing
thermal load
sensitivity}\label{a-non-linear-bayesian-hierarchical-model-for-characterizing-thermal-load-sensitivity}

\subsection{Model specification}\label{model-specification}

Assuming we have mortality data for different times and temperatures,
let \(y_{ij}\) denote the number of survivors out of \(n_{ij}\)
individuals in trial \(j\) exposed for time \(t_{ij}\) (e.g., minutes)
at assay temperature \(T_{ij}\) (°C). We can model the counts directly
with a binomial likelihood (but a beta-binomial likelihood function can
also be used and controls for overdispersion):

\begin{equation}\protect\phantomsection\label{eq-likelihood}{
y_{ij} \mid n_{ij}, p_{ij} \;\sim\; \text{Binomial}(n_{ij},\, p_{ij}),
}\end{equation}

with survival probability \(p_{ij}\) given by a four-parameter logistic
(4PL) function of \(\log_{10}\)-exposure time,

\begin{equation}\protect\phantomsection\label{eq-4pl}{
p_{ij} \;=\; \text{low} \;+\; \frac{\text{up} - \text{low}}{1 + \exp\!\bigl(k\,(\log_{10} t_{ij} - mid_{ij})\bigr)}.
}\end{equation}

where \(\text{low}\) is the lower survival probability asymptote,
\(\text{up}\) is the upper survival probability asymptote, and \(k\) is
the slope defining the change in survival probability on the
\(\log_{10} t\) axis. Larger \(k\) gives a sharper decline in survival
probability with \(\log_{10} t\). \(mid_{ij}\) is the value of
\(\log_{10} t\) at which \(p\) falls to the midpoint between
\(\text{low}\) and \(\text{up}\). When asymptotes are near 0 and 1,
\(mid_{ij}\) coincides with \(\log_{10}t_{50}\) at trial \(j\). While
other logistic forms can be used (e.g., two or 3 parameter logistics),
Equation~\ref{eq-4pl} is the most flexible model that can account for
different asymptotic survival probabilities.

The key parameters of Equation~\ref{eq-4pl} (i.e., \(\text{low}\),
\(\text{up}\), \(k\), \(mid_{ij}\)), can be fit as a non-linear function
within a likelihood or Bayesian framework, which estimates their
correlation and allows for fixed or random effects to be estimated to
explain variation in these parameters. Such an approach allows for
\emph{a priori} comparisons (e.g., do sexes or species differ in
midpoints?), with correct inference (e.g., do species differ
significantly from each other?) and allows for variance decomposition to
help identify key sources of variance in the data (e.g., estimating
variances explained across individuals, species, trials etc.). Such
models are most flexibly fit within a Bayesian framework with intuitive
interpretation of these parameters and a highly flexible path for error
propagation. We show below how such a model can be used to derive all
thermal load sensitivity parameters of interest to ecophysiologists.

\section{Deriving key thermal load sensitivity
parameters}\label{deriving-key-thermal-load-sensitivity-parameters}

The TLS framework is elegant in that it derives metrics with clear
biological interpretation that can be compared across populations and
species (e.g., \emph{z}, \(CT_{max_{1hr}}\)). Here, we show that we can
derived these same metrics directly from Equation~\ref{eq-4pl}.

As indicated above, parameters can be fit to model changes in the mean
of any of the four parameters in Equation~\ref{eq-4pl}. We can easily
derive thermal sensitivity (\emph{z}) and \(CT_{max_{1hr}}\) by
modelling the effect of temperature on midpoints as follows:

\begin{equation}\protect\phantomsection\label{eq-mid-temp}{
\text{mid}_{ij} = \beta_0 + \beta_1\,(T_{ij} - \bar{T}),
}\end{equation}

where \(\text{mid}_{ij}\) is the midpoint (\(\log_{10}\)-time at 50\%
survival between \(\text{low}\) and \(\text{up}\)) for trial \(j\) at
assay temperature \(T_{ij}\) (°C); \(\beta_0\) is the intercept giving
the midpoint at the grand mean temperature \(\bar{T}\); \(\beta_1\) is
the slope describing how the midpoint changes per °C increase in
\(T_{ij}\); and \(\bar{T}\) is the grand mean of assay temperatures used
to mean-centre \(T_{ij}\), which decorrelates \(\beta_0\) from
\(\beta_1\) and improves parameter estimation and interpretation
(\citeproc{ref-schielzeth_centering_2010}{Schielzeth 2010}).

Setting \(p_{ij} = 0.5\) in Equation~\ref{eq-4pl}, isolating the
exponential, and taking natural logs gives
\(\log_{10}t_{50}(T) = \text{mid}(T) + (1/k)\log[(\text{up} - 0.5)/(0.5 - \text{low})]\).
Substituting into Equation~\ref{eq-mid-temp} and grouping all terms that
do not depend on \(T\) into an intercept, \(\alpha\), recovers the
classical TDT linear relation:

\begin{equation}\protect\phantomsection\label{eq-t50}{
\log_{10} t_{50}(T) \;=\; \underbrace{\left[\beta_0 + \tfrac{1}{k}\log\!\tfrac{\text{up} - 0.5}{0.5 - \text{low}} - \beta_1\bar{T}\right]}_{\alpha} \;+\; \beta_1\,T,
}\end{equation}

Note that Equation~\ref{eq-t50} subtracts \(\beta_1\bar{T}\) in
\(\alpha\) because we mean-centred \(T\). With this subtraction, the
intercept estimates \(\log_{10}t_{50}\) at the average temperature
(\(\bar{T}\)); without it, the intercept estimates \(\log_{10}t_{50}\)
at \(T = 0\)°C. From Equation~\ref{eq-t50}, the classical TLS quantities
fall out directly. Thermal sensitivity is \(z = -1/\beta_1\), and the
critical thermal limit at any reference time \(t_\text{ref}\) is
\(CT_{max}(t_\text{ref}) = (\log_{10} t_\text{ref} - \alpha)/\beta_1\)
(with \(t_\text{ref} = 60\) min for \(CT_{max_{1hr}}\)). Each of these
quantities can be estimated separately for different strata (e.g., sex,
population, species) by including the relevant fixed covariates and/or
random effects in the linear predictor for \(\text{mid}_{ij}\) --- and,
if appropriate, in the linear predictors for \(\text{low}\),
\(\text{up}\), and \(k\), each of which is modelled separately in the
joint Bayesian fit (see Equation~\ref{eq-linpred} below).

The same derivation generalises beyond the median: the Bayesian
dose-response fit gives access to \emph{any} mortality percentile, not
just \(t_{50}\) (\citeproc{ref-jorgensen_unifying_2021}{Jørgensen
\emph{et al.} 2021};
\citeproc{ref-rezende_drosophila_mortality_2020}{Rezende \emph{et al.}
2020}). Setting \(p_{ij} = x/100\) in Equation~\ref{eq-4pl} and
following the same algebra used for \(t_{50}\), the posterior samples of
\((\text{low}, \text{up}, k, m(T))\) deliver the time at which \(x\%\)
of the population reaches the assayed endpoint:

\begin{equation}\protect\phantomsection\label{eq-LTx}{
\log_{10} t_x(T) \;=\; m(T) \;-\; \frac{1}{k}\,\ln\!\frac{\text{up} - x/100}{x/100 - \text{low}},
}\end{equation}

Varying \(x\) in Equation~\ref{eq-LTx} then yields any threshold of
interest. For instance, \(t_{10}\) marks the onset of early mortality
--- useful for risk management --- while \(t_{90}\) is closer to a
population-collapse threshold. As with \(t_{50}\), every prediction
inherits the joint posterior, so uncertainty bands can be reported
alongside the population-level median.

Importantly, the slope \(\beta_1\) --- and hence \(z\) --- is robust to
estimated asymptotes because the asymmetry correction
\(\tfrac{1}{k}\log\tfrac{\text{up} - 0.5}{0.5 - \text{low}}\) does not
depend on \(T\). Instead, it shifts \(\log_{10}t_{50}\) vertically
without changing the slope. In contrast, threshold summaries are not
robust in this way; when partial mortality or sparse data do not
estimate the asymptotes to be around 0 and 1, the asymmetry correction
must be retained --- added to the naive \(\log_{10}t_{50}\)
(Equation~\ref{eq-t50-corrected}) and subtracted, after division by the
slope (to convert the vertical shift in \(\log_{10}t_{50}\) into the
corresponding shift in temperature), from the naive \(CT_{max}\)
(Equation~\ref{eq-ctmax-corrected}):

\begin{equation}\protect\phantomsection\label{eq-t50-corrected}{
\log_{10}t_{50}(T) \;=\; \underbrace{\beta_0 + \beta_1(T - \bar{T})}_{\text{naive estimate}} \;+\; \underbrace{\tfrac{1}{k}\log\!\tfrac{\text{up} - 0.5}{0.5 - \text{low}}}_{\text{correction}},
}\end{equation}

\begin{equation}\protect\phantomsection\label{eq-ctmax-corrected}{
CT_{max}(t_\text{ref}) \;=\; \underbrace{\frac{\log_{10}t_\text{ref} - \beta_0 + \beta_1\bar{T}}{\beta_1}}_{\text{naive estimate}} \;-\; \underbrace{\tfrac{1}{\beta_1}\cdot\tfrac{1}{k}\log\!\tfrac{\text{up} - 0.5}{0.5 - \text{low}}}_{\text{correction}}.
}\end{equation}

In practice, however, these corrections may not be needed so long as the
experimental design includes treatments that span the full mortality
range which anchors \(\text{low}\) and \(\text{up}\) tightly near 1 and
0 and makes the correction term effectively zero. When this is the case,
\(\log_{10}t_{50}\) and \(CT_{max}\) can be read directly from the naive
forms in Equation~\ref{eq-t50-corrected} and
Equation~\ref{eq-ctmax-corrected}. Regardless, we recommend using the
correction factor as a safeguard. However, computing both the naive and
corrected estimates and comparing them is a useful sensitivity check for
whether the correction is needed in any given dataset.

\section{Predicting heat injury and survival under dynamic thermal
environments}\label{predicting-heat-injury-and-survival-under-dynamic-thermal-environments}

A particularly powerful aspect of TLS models are their ability to
translate dynamic temperature times series' from nature to predict heat
injury (HI) accumulation and thus population survival probability
(\citeproc{ref-jorgensen_unifying_2021}{Jørgensen \emph{et al.} 2021};
\citeproc{ref-li2023}{Li \emph{et al.} 2023};
\citeproc{ref-orsted_finding_2022}{Ørsted \emph{et al.} 2022},
\citeproc{ref-orsted_suzukii_2024}{2024};
\citeproc{ref-rezende_drosophila_mortality_2020}{Rezende \emph{et al.}
2020}). Heat injury accumulates as the time spent at a given temperature
divided by the median lethal time at that temperature, integrated across
the temperature trajectory \(T(i)\) as follows:

\begin{equation}\protect\phantomsection\label{eq-hi}{
HI(t) \;=\; 100 \int_0^t 10^{\,\big(T(i)\,-\,CT_{max_{1\text{hr}}}\big)/z} \,di,
}\end{equation}

where \(z\) is the thermal sensitivity and \(CT_{max_{1\text{hr}}}\) is
the static temperature at which 50\% of the population reaches the
assayed endpoint after one hour of exposure. When \(HI(t) = 100\%\), the
cumulative dose equals one \(t_{50}\) at the 1-hour reference,
predicting 50\% mortality (or 50\% loss of the assayed sublethal
endpoint, e.g.~fertility; Ørsted \emph{et al.}
(\citeproc{ref-orsted_suzukii_2024}{2024})). Injury accumulates only
when \(T(i)\) exceeds a critical temperature \(T_c\) below which damage
is repaired as fast as it accrues
(\citeproc{ref-jorgensen_unifying_2021}{Jørgensen \emph{et al.} 2021};
\citeproc{ref-orsted_finding_2022}{Ørsted \emph{et al.} 2022}). Field
temperature records are typically available at hourly resolution.
Equation~\ref{eq-hi} can then be evaluated by summing
\(100 \cdot 10^{(T_i - CT_{max_{1\text{hr}}})/z}\) across each hourly
interval whose temperature \(T_i\) exceeds \(T_c\), following the
implementations in Jørgensen \emph{et al.}
(\citeproc{ref-jorgensen_unifying_2021}{2021}), Ørsted \emph{et al.}
(\citeproc{ref-orsted_suzukii_2024}{2024}), and Li \emph{et al.}
(\citeproc{ref-li2023}{2023}). To propagate uncertainty, \(z\) and
\(CT_{max_{1\text{hr}}}\), calculated from our Bayesian dose-response
model can be used to calculate \(HI(t)\) through time by substituting
samples from the posterior distribution for each parameter into
Equation~\ref{eq-hi}.

A more natural prediction to make may instead be the population-level
survival probability over time. This has the advantage of allowing for a
direct connection to population dynamical models which typically use
survival, growth and fecundity as vital rates for population-level
predictions (\citeproc{ref-noble2026}{Noble \emph{et al.} 2026}). To
predict survival, we can calculate the \emph{effective exposure time} at
a chosen reference temperature \(T_{\text{ref}}\). More effective time
is accumulated at hot temperatures and less at cool ones, in proportion
to the scaling factor \(10^{(T(i) - T_{\text{ref}})/z}\),

\begin{equation}\protect\phantomsection\label{eq-tau-eff}{
\tau_{\text{eff}}(t) \;=\; \int_0^t 10^{\,(T(i) - T_{\text{ref}})/z}\,di,
}\end{equation}

We can then evaluate the 4PL surface once at the cumulative effective
exposure to obtain the predicted population-level survival fraction at
every point in the trajectory:

\begin{equation}\protect\phantomsection\label{eq-S-pred}{
S_{\text{pred}}(t) \;=\; S\!\big(T_{\text{ref}},\, \tau_{\text{eff}}(t)\big).
}\end{equation}

Here \(S(T, \tau)\) is the survival surface fitted by the dose-response
model, so \(S_{\text{pred}}(t)\) is the survival probability for an
organism held at \emph{constant} \(T_{\text{ref}}\) after exposure
duration \(\tau_{\text{eff}}(t)\).

\section{Application of Bayesian dose-response models to the TLS
framework}\label{application-of-bayesian-dose-response-models-to-the-tls-framework}

We used a four-parameter logistic function to model the impact of
temperature and duration (hours/minutes) on survival probability thermal
sensitivity, estimate \(CT_{max1h}\)

To break the well-known likelihood ridges among \((\text{low}, u, k)\)
in unconstrained 4PL fits
(\citeproc{ref-bolker_ecologicalmodels_2008}{Bolker 2008};
\citeproc{ref-gelman_priors_2006}{Gelman 2006}), we estimate the
identifiable parameterisation

\begin{equation}\protect\phantomsection\label{eq-reparam}{
\text{low} = 0.49\,\text{logit}^{-1}(\widetilde{\text{low}}), \qquad
\text{up}   = 0.51 + 0.49\,\text{logit}^{-1}(\widetilde{\text{up}}), \qquad
k   = \exp(\kappa),
}\end{equation}

with
\(\widetilde{\text{low}},\widetilde{\text{up}},\kappa \in \mathbb{R}\).
This constrains \(\text{low} \in [0, 0.49]\) and
\(\text{up} \in [0.51, 1]\) into disjoint intervals so that
\(u > \text{low}\) by construction, the curve direction is unambiguous,
and \(k\) is positive on the natural scale.

The inflection point \(m_{ij}\) is modelled as a linear function of
(mean-centred) assay temperature, with additional fixed covariates and
nested random effects:

\begin{equation}\protect\phantomsection\label{eq-linpred}{
m_{ij} \;=\; \beta_0 \;+\; \beta_1\,(T_{ij} - \bar T) \;+\; \mathbf{x}_{ij}^{\top}\boldsymbol{\gamma} \;+\; b_{j} \;+\; \varepsilon_{ij},
}\end{equation}

where \(\bar T\) is the grand mean of assay temperatures (centring
stabilises sampling and decorrelates \(\beta_0\) from \(\beta_1\)),
\(\mathbf{x}_{ij}\) holds any additional fixed covariates with
coefficients \(\boldsymbol{\gamma}\), \(b_{j}\) collects the nested
random effects and
\(\varepsilon_{ij} \sim \mathcal{N}(0, \sigma_\varepsilon^2)\) is an
observation-level random effect that absorbs the binomial overdispersion
characteristic of replicated count assays
(\citeproc{ref-harrison_olre_2014}{Harrison 2014}). However, we could
also remove the observation-level random effect from the model and
instead fit a beta-binomial error distribution which estimates a
dispersion parameter (\(\phi\)) more naturally.

Together, Equation~\ref{eq-likelihood}--Equation~\ref{eq-linpred} define
the hierarchical Bayesian 4PL fit directly to the raw binomial counts.
Every classical TDT quantity --- slope, \(z\), \(CT_{max}\) at any
reference time, and the heat-injury integral --- is recovered as a
one-line posterior summary of this single joint fit, as derived in
§{[}Equivalence with the classical pipeline{]}.

\section{Case Study 1: Thermal sensitivity estimation across life stages
in
shrimp}\label{case-study-1-thermal-sensitivity-estimation-across-life-stages-in-shrimp}

\subsection{What the 4PL delivers that the two-step
cannot}\label{what-the-4pl-delivers-that-the-two-step-cannot}

\section{Worked example: brown shrimp thermal
tolerance}\label{worked-example-brown-shrimp-thermal-tolerance}

\subsection{Data and experimental
design}\label{data-and-experimental-design}

\subsection{Lethal TDT --- hierarchical 4PL on
counts}\label{lethal-tdt-hierarchical-4pl-on-counts}

\subsection{Sublethal time-to-death}\label{sublethal-time-to-death}

\subsection{\texorpdfstring{Dynamic \(CT_{max}\) across ramping
rates}{Dynamic CT\_\{max\} across ramping rates}}\label{dynamic-ct_max-across-ramping-rates}

\subsection{Reconciling the three endpoints via the Jørgensen ramping
equation}\label{reconciling-the-three-endpoints-via-the-juxf8rgensen-ramping-equation}

\subsection{Field-temperature heat-injury projections under warming
scenarios}\label{field-temperature-heat-injury-projections-under-warming-scenarios}

\section{Discussion}\label{discussion}

\subsection{What changes for the field if the hierarchical 4PL becomes
the
default}\label{what-changes-for-the-field-if-the-hierarchical-4pl-becomes-the-default}

\subsection{Limits of the present
analysis}\label{limits-of-the-present-analysis}

\subsection{Recommendations for primary studies and
meta-analyses}\label{recommendations-for-primary-studies-and-meta-analyses}

\section{Conclusions}\label{conclusions}

\section{Acknowledgements}\label{acknowledgements}

\section*{Author contributions}\label{author-contributions}
\addcontentsline{toc}{section}{Author contributions}

\section*{Data and code availability}\label{data-and-code-availability}
\addcontentsline{toc}{section}{Data and code availability}

All data, analysis scripts, and the rendered manuscript and supplement
are available at
https://github.com/daniel1noble/tls\_model\_equivalence.

\section*{Conflict of interest}\label{conflict-of-interest}
\addcontentsline{toc}{section}{Conflict of interest}

The authors declare no competing interests.

\section*{References}\label{references}
\addcontentsline{toc}{section}{References}

\protect\phantomsection\label{refs}
\begin{CSLReferences}{1}{0}
\bibitem[\citeproctext]{ref-arnold_thermal_2025}
Arnold, P.A., Noble, D.W.A., Nicotra, A.B., Kearney, M.R., Rezende,
E.L., Andrew, S.C., \emph{et al.} (2025).
\href{https://doi.org/10.1111/gcb.70315}{A framework for modelling
thermal load sensitivity across life}. \emph{Global Change Biology}, 31,
e70315.

\bibitem[\citeproctext]{ref-bolker_ecologicalmodels_2008}
Bolker, B.M. (2008). \emph{Ecological models and data in {R}}. Princeton
University Press, Princeton, NJ.

\bibitem[\citeproctext]{ref-buckley_damage_2025}
Buckley, L.B., Huey, R.B. \& Ma, C.-S. (2025).
\href{https://doi.org/10.1093/icb/icaf019}{How damage, recovery, and
repair alter the fitness impacts of thermal stress}. \emph{Integrative
and Comparative Biology}, 65, 1061--1075.

\bibitem[\citeproctext]{ref-burke_ipd_2017}
Burke, D.L., Ensor, J. \& Riley, R.D. (2017).
\href{https://doi.org/10.1002/sim.7141}{Meta-analysis using individual
participant data: One-stage and two-stage approaches, and why they may
differ}. \emph{Statistics in Medicine}, 36, 855--875.

\bibitem[\citeproctext]{ref-carroll_measurement_2006}
Carroll, R.J., Ruppert, D., Stefanski, L.A. \& Crainiceanu, C.M. (2006).
\emph{\href{https://doi.org/10.1201/9781420010138}{Measurement error in
nonlinear models: A modern perspective}}. 2nd edn. {Chapman and
Hall/CRC}.

\bibitem[\citeproctext]{ref-deutsch_latitude_2008}
Deutsch, C.A., Tewksbury, J.J., Huey, R.B., Sheldon, K.S., Ghalambor,
C.K., Haak, D.C., \emph{et al.} (2008).
\href{https://doi.org/10.1073/pnas.0709472105}{Impacts of climate
warming on terrestrial ectotherms across latitude}. \emph{Proceedings of
the National Academy of Sciences}, 105, 6668--6672.

\bibitem[\citeproctext]{ref-gelman_priors_2006}
Gelman, A. (2006). \href{https://doi.org/10.1214/06-BA117A}{Prior
distributions for variance parameters in hierarchical models (comment on
article by {Browne} and {Draper})}. \emph{Bayesian Analysis}, 1,
515--534.

\bibitem[\citeproctext]{ref-harrison_olre_2014}
Harrison, X.A. (2014). \href{https://doi.org/10.7717/peerj.616}{Using
observation-level random effects to model overdispersion in count data
in ecology and evolution}. \emph{PeerJ}, 2, e616.

\bibitem[\citeproctext]{ref-harrison_mixed_2018}
Harrison, X.A., Donaldson, L., Correa-Cano, M.E., Evans, J., Fisher,
D.N., Goodwin, C.E.D., \emph{et al.} (2018).
\href{https://doi.org/10.7717/peerj.4794}{A brief introduction to mixed
effects modelling and multi-model inference in ecology}. \emph{PeerJ},
6, e4794.

\bibitem[\citeproctext]{ref-jorgensen_drosophila_2019}
Jørgensen, L.B., Malte, H. \& Overgaard, J. (2019).
\href{https://doi.org/10.1111/1365-2435.13279}{How to assess
{Drosophila} heat tolerance: Unifying static and dynamic tolerance
assays to predict heat distribution limits}. \emph{Functional Ecology},
33, 629--642.

\bibitem[\citeproctext]{ref-jorgensen_unifying_2021}
Jørgensen, L.B., Malte, H. \& Overgaard, J. (2021).
\href{https://doi.org/10.1038/s41598-021-92004-6}{A unifying model to
estimate thermal tolerance limits in ectotherms across static, dynamic
and fluctuating exposures to thermal stress}. \emph{Scientific Reports},
11, 12840.

\bibitem[\citeproctext]{ref-jorgensen_nature_2022}
Jørgensen, L.B., Ørsted, M., Malte, H., Wang, T. \& Overgaard, J.
(2022). \href{https://doi.org/10.1038/s41586-022-05334-4}{Extreme
escalation of heat failure rates in ectotherms with global warming}.
\emph{Nature}, 611, 93--98.

\bibitem[\citeproctext]{ref-li2023}
Li, Y.-J., Chen, S.-Y., Jørgensen, L.B., Overgaard, J., Renault, D.,
Colinet, H., \emph{et al.} (2023).
\href{https://doi.org/10.1016/j.jtherbio.2023.103583}{Interspecific
differences in thermal tolerance landscape explain aphid community
abundance under climate change}. \emph{Journal of Thermal Biology}, 114,
103583.

\bibitem[\citeproctext]{ref-nakagawa_santos_2012}
Nakagawa, S. \& Santos, E.S.A. (2012).
\href{https://doi.org/10.1007/s10682-012-9555-5}{Methodological issues
and advances in biological meta-analysis}. \emph{Evolutionary Ecology},
26, 1253--1274.

\bibitem[\citeproctext]{ref-nielsen_nlmm_2004}
Nielsen, O.K., Ritz, C. \& Streibig, J.C. (2004).
\href{https://doi.org/10.1614/WT-03-070R1}{Nonlinear mixed-model
regression to analyze herbicide dose-response relationships}. \emph{Weed
Technology}, 18, 30--37.

\bibitem[\citeproctext]{ref-noble2026}
Noble, D.W.A., Mayfield, M.M., Hoffmann, A.A., Chen, Z.-H., Lade, S.J.,
Bai, X., \emph{et al.} (2026).
\href{https://doi.org/10.1016/j.tree.2026.01.009}{A systems modelling
approach to predict biological responses to extreme heat}. \emph{Trends
in Ecology \& Evolution}, 41, 451--464.

\bibitem[\citeproctext]{ref-orsted_finding_2022}
Ørsted, M., Jørgensen, L.B. \& Overgaard, J. (2022).
\href{https://doi.org/10.1242/jeb.244514}{Finding the right thermal
limit: A framework to reconcile ecological, physiological and
methodological aspects of {CTmax} in ectotherms}. \emph{Journal of
Experimental Biology}, 225, jeb244514.

\bibitem[\citeproctext]{ref-orsted_suzukii_2024}
Ørsted, M., Willot, Q., Olsen, A.K., Kongsgaard, V. \& Overgaard, J.
(2024). \href{https://doi.org/10.1111/ele.14421}{Thermal limits of
survival and reproduction depend on stress duration: A case study of
{Drosophila} suzukii}. \emph{Ecology Letters}, 27, e14421.

\bibitem[\citeproctext]{ref-pinsky_marine_2019}
Pinsky, M.L., Eikeset, A.M., McCauley, D.J., Payne, J.L. \& Sunday, J.M.
(2019). \href{https://doi.org/10.1038/s41586-019-1132-4}{Greater
vulnerability to warming of marine versus terrestrial ectotherms}.
\emph{Nature}, 569, 108--111.

\bibitem[\citeproctext]{ref-pottier_amphibians_2025}
Pottier, P., Kearney, M.R., Wu, N.C., Gunderson, A.R., Rej, J.E.,
Rivera-Villanueva, A.N., \emph{et al.} (2025).
\href{https://doi.org/10.1038/s41586-025-08665-0}{Vulnerability of
amphibians to global warming}. \emph{Nature}, 639, 954--961.

\bibitem[\citeproctext]{ref-rezende_drosophila_mortality_2020}
Rezende, E.L., Bozinovic, F., Szilágyi, A. \& Santos, M. (2020).
\href{https://doi.org/10.1126/science.aba9287}{Predicting temperature
mortality and selection in natural {Drosophila} populations}.
\emph{Science}, 369, 1242--1245.

\bibitem[\citeproctext]{ref-rezende_landscapes_2014}
Rezende, E.L., Castañeda, L.E. \& Santos, M. (2014).
\href{https://doi.org/10.1111/1365-2435.12268}{Tolerance landscapes in
thermal ecology}. \emph{Functional Ecology}, 28, 799--809.

\bibitem[\citeproctext]{ref-ritz_ecotox_2010}
Ritz, C. (2010). \href{https://doi.org/10.1002/etc.7}{Toward a unified
approach to dose-response modeling in ecotoxicology}.
\emph{Environmental Toxicology and Chemistry}, 29, 220--229.

\bibitem[\citeproctext]{ref-ritz_drc_2015}
Ritz, C., Baty, F., Streibig, J.C. \& Gerhard, D. (2015).
\href{https://doi.org/10.1371/journal.pone.0146021}{Dose-response
analysis using {R}}. \emph{PLoS ONE}, 10, e0146021.

\bibitem[\citeproctext]{ref-ritz_streibig_2008}
Ritz, C. \& Streibig, J.C. (2008).
\emph{\href{https://doi.org/10.1007/978-0-387-09616-2}{Nonlinear
regression with {R}}}. Springer-Verlag, New York.

\bibitem[\citeproctext]{ref-rudemo_random_1989}
Rudemo, M., Ruppert, D. \& Streibig, J.C. (1989).
\href{https://doi.org/10.2307/2531483}{Random-effect models in nonlinear
regression with applications to bioassay}. \emph{Biometrics. Journal of
the International Biometric Society}, 45, 349--362.

\bibitem[\citeproctext]{ref-schielzeth_centering_2010}
Schielzeth, H. (2010).
\href{https://doi.org/10.1111/j.2041-210X.2010.00012.x}{Simple means to
improve the interpretability of regression coefficients}. \emph{Methods
in Ecology and Evolution}, 1, 103--113.

\bibitem[\citeproctext]{ref-smale_marine_2019}
Smale, D.A., Wernberg, T., Oliver, E.C.J., Thomsen, M., Harvey, B.P.,
Straub, S.C., \emph{et al.} (2019).
\href{https://doi.org/10.1038/s41558-019-0412-1}{Marine heatwaves
threaten global biodiversity and the provision of ecosystem services}.
\emph{Nature Climate Change}, 9, 306--312.

\bibitem[\citeproctext]{ref-snook_fertility_2026}
Snook, R.R., Bretman, A., Dougherty, L.R. \& Fricke, C. (2026).
\href{https://doi.org/10.1038/s44358-026-00142-4}{The consequences of
rising temperatures for animal fertility}. \emph{Nature Reviews
Biodiversity}, 2, 225--241.

\bibitem[\citeproctext]{ref-sunday_global_2011}
Sunday, J.M., Bates, A.E. \& Dulvy, N.K. (2011).
\href{https://doi.org/10.1098/rspb.2010.1295}{Global analysis of thermal
tolerance and latitude in ectotherms}. \emph{Proceedings of the Royal
Society B: Biological Sciences}, 278, 1823--1830.

\bibitem[\citeproctext]{ref-sunday_redistribution_2012}
Sunday, J.M., Bates, A.E. \& Dulvy, N.K. (2012).
\href{https://doi.org/10.1038/nclimate1539}{Thermal tolerance and the
global redistribution of animals}. \emph{Nature Climate Change}, 2,
686--690.

\bibitem[\citeproctext]{ref-truebano2018}
Truebano, M., Fenner, P., Tills, O., Rundle, S. \& Rezende, E. (2018).
\href{https://doi.org/10.1242/jeb.171629}{Thermal strategies vary with
life history stage}. \emph{Journal of Experimental Biology}, 221,
jeb171629.

\bibitem[\citeproctext]{ref-villeneuve_dynamic_2024}
Villeneuve, A.R. \& White, E.R. (2024).
\href{https://doi.org/10.1111/1365-2656.14120}{Predicting organismal
response to marine heatwaves using dynamic thermal tolerance landscape
models}. \emph{Journal of Animal Ecology}, 94, 1515--1527.

\bibitem[\citeproctext]{ref-wiens_localextinctions_2016}
Wiens, J.J. (2016).
\href{https://doi.org/10.1371/journal.pbio.2001104}{Climate-related
local extinctions are already widespread among plant and animal
species}. \emph{PLOS Biology}, 14, e2001104.

\end{CSLReferences}




\end{document}
